\documentclass[12pt,a4paper]{article}
\usepackage[T1]{fontenc}
\title{Errata and Questions}
\author{Daniel Martling}
\date{}
%--Packages--
\usepackage[english]{babel}
\usepackage[utf8]{inputenc}
\usepackage[T1]{fontenc}
%\usepackage{showkeys}

\usepackage{latexsym,amsmath,amsthm,amssymb,bm,lmodern,titlesec,xcolor,graphicx,setspace,pdfpages,
	geometry,marginnote,caption,array,sectsty,siunitx,eurosym,colortbl,float,wrapfig,vmargin,tikz,hyperref}
\usepackage[nottoc]{tocbibind}
\usepackage{tikz-cd}
%\usepackage{}
\usepackage{lipsum}
\usepackage{url}

\setpapersize{A4}
\setmarginsrb{30mm}{20mm}{30mm}{30mm}%left top right bottom
{8mm}{6mm}{10mm}{10mm}%headheight headsep footheight footskip

% Some useful colors
\definecolor{Rosso}{RGB}{232,61,61}
\definecolor{Viola}{RGB}{188,30,188}
\definecolor{Celeste}{RGB}{34,186,211}
\definecolor{Arancione}{RGB}{232,181,19}
\definecolor{Blu}{RGB}{0, 44, 92}
\definecolor{Azzurro}{RGB}{57,186,238}


\hypersetup{colorlinks=true,linkcolor=Blu,allbordercolors=white, urlcolor=blue, citecolor=Viola}
\usepackage{enumitem} 
\renewcommand\labelitemi{\textbullet} 


% math environments
%  we use \newtheorem to specify our theorem-like environments
% (theorem, definition, etc.) and how to display and number them.
%
% Note: The \theoremstyle declarations affect the appearance of the
% Theorems, Definitions, etc.
% Note 2: If you want swedish names you have to go and change the argument in brackets e.g. \newtheorem{thm}[defi]{Theorem} should become \newtheorem{thm}[defi]{Sats}
\theoremstyle{definition}
\newtheorem{definition}{Definition}[section]

%\newtheorem{hottake}{Hot Take}[section]

\newtheorem*{problem}{Problem}
\newtheorem{conjecture}[definition]{Conjecture}

%\theoremstyle{plain}
\newtheorem{theorem}[definition]{Theorem}
\newtheorem{proposition}[definition]{Proposition}
\newtheorem{corollary}[definition]{Corollary}
\newtheorem{lemma}[definition]{Lemma}
%\newtheorem*{remark}{Remark}%[definition]{Remark}
%\newtheorem{example}[definition]{Example}
%\newtheorem*{notation}{Notation}
%\newtheorem*{note}{Note}

\theoremstyle{remark}
\newtheorem*{question}{Question}
\newtheorem{remark}[definition]{Remark}
\newtheorem{example}[definition]{Example}
\newtheorem*{notation}{Notation}
\newtheorem*{note}{Note}

%\onehalfspacing
\numberwithin{equation}{section}

%--Heading of  section--
\titleformat{\section}[display]
{\normalfont\bfseries\color{Blu}}
{\filleft
	\begin{tikzpicture}
		\node[
		outer sep=0pt,
		text width=2.5cm,
		minimum height=3cm,
		fill=white,
		font=\color{Blu}\fontsize{80}{0}\selectfont,
		align=center
		] (num) {\thesection};
		\node[
		rotate=90,
		anchor=south,
		font=\color{black}\large\normalfont
		] at ([xshift=-5pt]num.west) {\Huge\textsc{}};  
	\end{tikzpicture}%
}
{10pt}
{\titlerule\vskip4pt\Huge\sc\sffamily}

%--Heading of subsection--
\newcommand*\numb[1]{%
	\begin{tikzpicture}[baseline=-0.7ex]
		\node[
		outer sep=0pt,
		text width=0.6cm,
		minimum height=0.6cm,
		fill=Blu,
		font=\color{white}\fontsize{12}{20}\selectfont,
		align=center
		] (num) {\thesubsection};
	\end{tikzpicture}%
}
\titleformat{\subsection}
{\normalfont\color{Blu}\large\sc\sffamily}{\numb{\thesubsection}}{0.8em}{}

%--Heading subsubsection--
\newcommand*\stocaz[1]{%
	\begin{tikzpicture}[baseline=-0.7ex]
		\node[
		outer sep=0pt,
		text width=0.8cm,
		minimum height=0.5cm,
		fill=Blu,
		font=\color{white}\fontsize{10}{20}\selectfont,
		align=center
		] (num) {\thesubsubsection};
	\end{tikzpicture}%
}
\titleformat{\subsubsection}
{\normalfont\color{Blu}\bfseries\sc\sffamily}{\stocaz{\thesubsubsection}}{0.8em}{}

% Blackboard Boldface
\newcommand{\NN}{\mathbb{N}}
\newcommand{\ZZ}{\mathbb{Z}}
\newcommand{\QQ}{\mathbb{Q}}
\newcommand{\RR}{\mathbb{R}}
\newcommand{\CC}{\mathbb{C}}
\newcommand{\KK}{\mathbb{K}}
\newcommand{\FF}{\mathbb{F}}

\newcommand{\1}{\mathbf{1}}
\newcommand{\0}{\mathbf{0}}

% Text
\newcommand{\GL}{\text{GL}}
\newcommand{\SL}{\text{SL}}
\newcommand{\SymP}{\text{Sym}}
\newcommand{\AltP}{{\bigwedge}}
\newcommand{\Hom}{\text{Hom}}
\newcommand{\End}{\text{End}}
\newcommand{\SO}{\text{SO}}
\newcommand{\SU}{\text{SU}}
\newcommand{\id}{\text{id}}
\newcommand{\Stab}{\text{Stab}}
\newcommand{\Cent}{\text{Cent}}
\newcommand{\Norm}{\text{Norm}}
\newcommand{\Type}{\text{Type}}
\newcommand{\Dih}{\text{D}}
\newcommand{\Sym}{\mathcal{S}} %{\mathfrak{S}} % \text{S}
%\newcommand{\Alt}{\text{A}}
\newcommand{\Klein}{\text{V}}
\newcommand{\Cyc}{\mathcal{C}}%{\text{C}}
\newcommand{\sgn}{\text{sgn}}
\newcommand{\Span}{\text{Span}}
\newcommand{\im}{\text{im }}
\newcommand{\Tr}{\text{Tr }}
\newcommand{\Reg}{\text{Reg}}
\newcommand{\Triv}{\text{Triv}}
\newcommand{\Alt}{\text{Alt}}
\newcommand{\Perm}{\text{Perm}}
\newcommand{\Stan}{\text{Stan}}
% Vectors
\newcommand{\bas}{\mathbf{\hat{e}}}
\newcommand{\fbas}{\mathbf{\hat{f}}}

\newcommand{\ihat}{\mathbf{\hat{\imath}}}
\newcommand{\jhat}{\mathbf{\hat{\jmath}}}
\newcommand{\ghat}{\mathbf{\hat{g}}}
\newcommand{\hhat}{\mathbf{\hat{h}}}
\newcommand{\vhat}{\mathbf{\hat{v}}}
\newcommand{\what}{\mathbf{\hat{w}}}
\newcommand{\uhat}{\mathbf{\hat{u}}}
\newcommand{\xhat}{\mathbf{\hat{x}}}

\newcommand{\vvec}{\mathbf{v}}
\newcommand{\uvec}{\mathbf{u}}
\newcommand{\wvec}{\mathbf{w}}


\newcommand{\Ss}{\scriptstyle}

%\newcommand{\barpi}{\overline{\pi}}
\newcommand{\barpi}{\bar{\pi}}

\allowdisplaybreaks

\begin{document}
	\maketitle
	
	\section*{Errors courtesy of Håkan}
	
		\paragraph{Page 6, Sammanfattning.} Irreducible representation is translated to ``irreducibla representationer'' in Swedish, not ``oreducerbara representationer''.
		
		\paragraph{Page 13, Example 2.9.} Property iii) of a projection should be ``$\pi(\vvec) = \0$ for any $\vvec$ in $V$ but not in $W$''.
		
	\section*{Other errors found}
	
		\paragraph{Page 19, Degree one repr. of $\Cyc_n$.} The map should be to the multiplicative group of complex numbers, that is $\rho: \Cyc_n \rightarrow \CC^\times$.
		
		\paragraph{Page 31, Proof of Corollary 3.26.} $h$ and $g$ should be exchanged in the left hand side of the calculation, that is it should be \begin{align*}
			\rho_g \cdot \rho_h &= \rho_{gh} \\
			&= \rho_{hg} \\
			&= \rho_h \cdot \rho_g.
		\end{align*}
		
		\paragraph{Page 38, Inner product property iii).} The expression for $(\varphi|\varphi)$ holds for \textit{any} character $\varphi$, however in particular for non-zero characters we have that $(\varphi|\varphi) > 0$ and $(\varphi|\varphi) \in \RR$.
		
		\paragraph{Page 40, Example 4.17.} The calculation for the standard character should be 
		\begin{align*}
			(\chi_\Stan|\chi_\Stan) &= \frac{1}{24}\left( 3^2 + 6 \cdot 1^2 + 8 \cdot 0^2 + 6 \cdot (-1)^2 + 3 \cdot (-1)^2 \right) \\
			&= 1.
		\end{align*}
		
		\paragraph{Page 41, Multiplicities in regular representation decomposition.} The summation sign in the comment on the seventh row of the calculation for $m_j$ should be a direct sum sign, that is \begin{align*}
			V = \bigoplus_i m_i V_i.
		\end{align*}
		
		\paragraph{Page 44, Decomposition of $V_\Stan \otimes V_\Stan$.} The calculation for $(\chi_\Alt|\chi_\Stan^2)$ should be \begin{align*}
			\frac{1}{24}\left(9-6+0-6+{\emph{3}}\right) = 0.
		\end{align*}
		
	\section*{Questions from Håkan}\text
	
	\begin{enumerate}
		\item[1.] På s.25, exempel 3.17: Vad är den precisa ekvationen för $\tau(\wvec)$ som du borde visat här?
			
			\textbf{Answer:} Let $G = \Sym_n$ and $V$ be the regular representation space spanned by the basis $\{\bas_\sigma\}_{\sigma \in \Sym_n}$. Let $W$ be the subspace of $V$ spanned by the sum
			\begin{align*}
				\sum_{\sigma \in \Sym_n} \sgn(\sigma) \bas_\sigma.
			\end{align*}
			A vector $\wvec \in W$ is a scalar multiple of this sum, that is for some $a \in \CC$,
			\begin{align*}
				\wvec = a \sum_{\sigma \in \Sym_n} \sgn(\sigma) \bas_\sigma.
			\end{align*}
			The action of a $\tau \in \Sym_n$ in this $\wvec$ is
			\begin{align*}
				\tau(\wvec) &= \tau \cdot a \sum_{\sigma \in \Sym_n} \sgn(\sigma) \bas_\sigma \\
				&= a \sum_{\sigma \in \Sym_n} \sgn(\sigma) \tau \cdot \bas_\sigma & \text{($\tau$ is linear)} \\
				&= a \sum_{\sigma \in \Sym_n} \sgn(\sigma) \cdot \bas_{\tau\sigma} & \text{(Action of $\Sym_n$ on basis)} \\
				&= a \frac{1}{\sgn(\tau)} \sum_{\sigma \in \Sym_n} \sgn(\tau)\sgn(\sigma) \cdot \bas_{\tau\sigma} \\
				&= a \frac{1}{\sgn(\tau)} \sum_{\sigma \in \Sym_n} \sgn(\tau\sigma) \cdot \bas_{\tau\sigma} & \text{($\sgn$ is a homom.)}
			\end{align*}
			Now, let $\sigma' = \tau\sigma$, then we can let the sum be indexed over these $\sigma'$ since when $\sigma$ goes over all of $\Sym_n$, so will $\sigma'$. Also since $\sgn(\tau) = \pm 1$ we have that $\sgn(\tau)^{-1} = \sgn(\tau)$ hence
			\begin{align*}
				\tau(\wvec) &= \sgn(\tau) \cdot a \sum_{\sigma' \in \Sym_n} \sgn(\sigma') \bas_{\sigma'} \\
				&= \sgn(\tau)\wvec.
			\end{align*}
			We have found that $\tau$ changes any vector in $W$ with a factor $\pm 1$ hence $W$ is a $\Sym_n$-invariant subspace of $V$.
			
		\item[2.] Vad kan man säga om totala antalet irreducibla representationer för en ändlig grupp $G$?
		
			 \textbf{Answer:} Firstly, as found in the text, the square of the dimensions of each irreducible representation will sum to the order of the group. However, for explicitly finding the expected number of irreducibles, this can be said:
			 
			 Analogous to a basis in a vector space, the irreducible characters of the group create an orthonormal system in a complex vector space of characters. In this analogy, any arbitrary character is a linear combination of irreducibles and the dimension of the vector space is the number of entries in the character which is the number of conjugacy classes, since the trace is a class function (=constant over classes). This is assuming there are no non-zero class functions perpendicular to this basis (Fulton Harris). It could then be said that there are as many irreducible representations of a finite group as there are conjugacy classes in that group.
			
			For example, in the abelian $\Cyc_n$ every element has its own conjugacy class, hence there are $n$ irreducible (degree one) representations of $\Cyc_n$. Another example is the symmetric groups where the number of irreducible representations of $\Sym_n$ are the number of integer partitions of $n$.
			
		\item[3.] Du visar att en abelsk grupp bara har 1-dimensionella representationer. Vad kan man säga i allmänhet om antalet 1-dimensionella representationer av en grupp $G$?
		
			 \textbf{Answer:} Låt $G$ vara en (icke-abelsk) grupp och låt $H$ vara en abelsk delgrupp till $G$. Då är alla irreducibla representationer till $H$ av grad 1. Antalet sådana representationer av $H$ är då $|H|$. 
			 
			 
			 
			 Commutator subgroup $[G,G]$ - subgroup of $G$ generated by all elements $g^{-1}h^{-1}gh$ (including all non-commuting elements basically). For abelian group, $[G,G] = \{e\}$. Number of degree one representations of $G$ is the size of the quotient group $G/[G,G]$. Every symmetric group $\Sym_n$ has a normal subgroup $\text{A}_n$ of index 2 isomorphic to $[\Sym_n, \Sym_n]$, hence the number of degree one irreducible representations of $\Sym_n$ is 2.
			
		\item[4.] Kan du bestämma karaktärstabellen för någon annan grupp, t ex för den alternerande gruppen $\text{A}_4$, utifrån de egenskaper vi vet karaktärerna har?
		
			\textbf{Answer:} För till exempel $\Sym_5$ behövs symmetriska och yttre kvadraterna av standardrepresentationen (två komplementerande delrum till standardrepresentationens tensorkvadrat) för att hitta alla irreducibla representationer.
			
			För $\text{A}_4$ kan det noteras att den är en normal delgrupp till $\Sym_4$ med index 2, och begränsas representationerna av $\Sym_4$ till $\text{A}_4$ ges att den triviala och standardrepresentationerna är irreducibla för $\text{A}_4$.
			
	\end{enumerate}
				
\end{document}